\chapter{归约、扫描与分区}

先看一个并行 filter 的常见事故。多个线程扫描输入，遇到满足条件的元素就把它 \code{push_back} 到同一个输出 vector。加锁后结果正确但很慢；改成原子递增下标后，输出位置唯一了，顺序却变成线程调度顺序；改成每个线程一个本地 vector 后，最后拼接又要说明按什么顺序拼接，失败后哪些元素有效也说不清。问题不在于没有找到“更快的容器”，而在于并行程序没有先回答：每个输出元素到底写到哪里。

本章从这个问题进入三个基础模式：\term{归约}{reduce}、\term{扫描}{scan} 和 \term{分区}{partition}。Reduce 把全局状态拆成局部 partial 再合并；scan 把“前面已经输出多少”变成可计算的前缀位置；partition 把多个输出目的地的写入区间提前分配好。它们不是算法题里的孤立名词，而是数据库执行器、日志计算、图处理、排序、shuffle、checkpoint 和分布式恢复共同使用的系统骨架。

\begin{learninggoals}
学完本章后，读者应能完成六件事：
\begin{enumerate}
  \item 从串行 reference 推导并行切分方式，而不是先开线程再补正确性。
  \item 为 reduce 写出 identity、合并函数、确定性、浮点误差和 partial 所有权。
  \item 解释 scan 为什么能把稳定 filter 的共享 \code{push_back} 改成无冲突写区间。
  \item 为多路 partition 推导二维计数、前缀偏移、写出区间和 manifest。
  \item 判断静态切分、动态调度、热点隔离、稳定输出和 cutoff 的取舍。
  \item 把单机 partition 的 block 身份迁移到分布式 shuffle 和失败重做。
\end{enumerate}
\end{learninggoals}

\begin{systemdiagram}{从串行隐含状态到并行显式结构}
\begin{pdfdiagram}
\node[book accent node,text width=2.55cm] (reduce) at (0,0) {串行累计值\\sum/count/state};
\node[book teal node,text width=2.55cm] (scan) at (3.35,0) {串行输出位置\\push count};
\node[book purple node,text width=2.55cm] (partition) at (6.7,0) {串行桶状态\\bucket append};

\node[book small node,text width=2.55cm] (partial) at (0,-1.55) {局部 partial\\固定合并树};
\node[book small node,text width=2.55cm] (offset) at (3.35,-1.55) {exclusive scan\\输出起点};
\node[book small node,text width=2.55cm] (manifest) at (6.7,-1.55) {bucket range\\block manifest};

\draw[book strong arrow] (reduce) -- (partial);
\draw[book strong arrow] (scan) -- (offset);
\draw[book strong arrow] (partition) -- (manifest);

\node[book label,text width=9.4cm] at (3.35,-2.75) {并行算法的关键不是线程 API，而是把串行循环里的隐含状态显式化为 partial、offset 和 manifest。};
\end{pdfdiagram}
\diagramnote{本图要证明的命题是：reduce、scan、partition 都是在把串行状态改造成可合并、可定位、可恢复的结构。}
\end{systemdiagram}

\section{从串行语义推导并行结构}

并行算法不能从“开几个线程”开始。第一步永远是串行语义：输入中哪些元素参与计算，输出顺序是否重要，错误如何报告，浮点误差如何比较，输出是否可重试。串行 reference 可以慢，但它定义了问题。并行版本的每一次重排、切分和合并，都必须能回到 reference 解释。

\topic{Reference 和合同}
以 \code{count_if} 为例，串行 reference 只做一件事：从左到右检查每个元素，满足谓词就加一。

\begin{lstlisting}[language=C++,caption={串行 count-if reference}]
std::int64_t count_at_least(std::span<const std::int32_t> values,
                            std::int32_t threshold) {
  std::int64_t count = 0;
  for (const std::int32_t value : values) {
    if (value >= threshold) {
      ++count;
    }
  }
  return count;
}
\end{lstlisting}

这个 reference 暗含四个事实：每个输入元素被检查一次且只检查一次；谓词没有外部副作用；输出只是一个整数；整数加法是合并规则。并行版本可以把输入切成块，是因为这些事实允许局部计数和最后求和。若问题改成“输出所有满足谓词的元素并保持顺序”，算法结构立刻改变，因为输出位置不再是单个计数器。

\topic{半开区间和覆盖证明}
数组切分应使用半开区间 \code{[begin, end)}。好的切分满足三个条件：所有块覆盖整个输入；相邻块不重叠；线程数大于元素数时仍然正确。一个常用公式如下：

\begin{lstlisting}[language=C++,caption={按块编号计算半开区间}]
struct Range {
  std::size_t begin = 0;
  std::size_t end = 0;
};

Range block_range(std::size_t total,
                  std::size_t block_index,
                  std::size_t block_count) {
  const std::size_t begin = (total * block_index) / block_count;
  const std::size_t end = (total * (block_index + 1)) / block_count;
  return Range{begin, end};
}
\end{lstlisting}

工程实现还要处理 \code{block_count == 0} 和极大输入乘法溢出。教材中先用这个公式，是为了让读者看到覆盖关系：第 \code{k} 块的 \code{end} 等于第 \code{k + 1} 块的 \code{begin}，第一块从 0 开始，最后一块到 \code{total} 结束。并行正确性的很多证明，都从这个切分证明开始。

\topic{ParallelSemanticContract}
本章建议在项目中引入一个教学用合同对象，记录并行化前必须回答的问题。

\begin{lstlisting}[language=C++,caption={ParallelSemanticContract 的教学形态}]
enum class OutputOrder {
  Stable,
  Unordered,
  DeterministicByKey
};

struct ParallelSemanticContract {
  bool has_identity = true;
  bool associative = true;
  bool commutative = true;
  bool allows_floating_error = false;
  OutputOrder output_order = OutputOrder::Stable;
  bool retryable_by_chunk = false;
};
\end{lstlisting}

这个结构不是要把所有算法抽象成一个框架，而是把思考顺序固定下来。Reduce 关注 identity 和合并函数；scan 关注输出位置；partition 关注目标桶和 manifest；恢复关注 chunk 是否可重做。若合同写不清，就不应进入线程实现。

\section{Reduce：把共享状态变成 partial}

Reduce 的第一原则是避免高频共享写。多个线程对一个全局原子计数器 \code{fetch_add}，结果可能正确，但所有 worker 都在争同一条 cache line。线程本地 partial 把高频更新留在本地，阶段末尾再合并。这个思想在 Linux per-cpu counter、数据库 map-side combine、分布式 reduce 和高性能 telemetry 中反复出现。

\topic{Identity、结合律和确定性}
Reduce 需要 identity 和合并函数。整数求和的 identity 是 0，最大值的 identity 可以是类型最小值，集合并的 identity 是空集合。若 identity 选错，空输入和空块就会错。结合律决定能否改变合并顺序；交换律决定能否重排 partial；确定性决定输出是否每次一致。

浮点求和是最好的反例。数学实数加法满足结合律，但机器浮点加法会舍入。串行左折叠、四累加器、树形合并、按 NUMA 节点分层合并，可能得到不同低位结果。若业务需要逐位可复现，就要固定合并树；若业务允许误差，就要写明绝对误差、相对误差或 ULP 边界。并行化之前必须选择语义。

\topic{Partial 的所有权}
线程本地 partial 应避免 false sharing。每个 worker 写自己的槽位，最后由合并阶段读取。若 partial 很小，可以用对齐槽；若 partial 是大型 histogram 或 hash map，内存预算和合并策略就成为核心问题。

\begin{lstlisting}[language=C++,caption={线程本地 partial 的基本形态}]
constexpr std::size_t CACHE_LINE_BYTES = 64;

struct alignas(CACHE_LINE_BYTES) CountPartial {
  std::int64_t count = 0;
};

void count_worker(std::span<const std::int32_t> values,
                  std::int32_t threshold,
                  CountPartial& partial) {
  std::int64_t local = 0;
  for (const std::int32_t value : values) {
    if (value >= threshold) {
      ++local;
    }
  }
  partial.count = local;
}
\end{lstlisting}

这个 worker 只写自己的 \code{partial}。正确性来自区间切分和整数加法；性能收益来自消除每个命中元素上的全局同步。代价也要写清：partial 占用额外空间，读取全局结果有延迟，最后合并可能成为新瓶颈。

\topic{线性合并和树形合并}
线性合并最简单：一个线程从头到尾扫描 partial。worker 数少、partial 小时足够。树形合并把 partial 两两合并，关键路径从 \complexity{O(p)} 降到 \complexity{O(\log p)}，其中 \code{p} 是 partial 数。树形合并的代价是更多阶段、更多临时对象和更复杂调度。

\begin{center}
\begin{tabular}{p{0.23\linewidth}p{0.31\linewidth}p{0.34\linewidth}}
\toprule
合并方式 & 适用场景 & 主要风险 \\
\midrule
线性合并 & partial 很小，worker 数少，简单优先 & 单线程尾部，合并阶段不可并行 \\
固定树形合并 & 需要可复现，partial 较多 & 调度自由度降低，临时对象更多 \\
动态树形合并 & 吞吐优先，partial 大小不均 & 浮点或顺序敏感结果可能不稳定 \\
分层合并 & NUMA 或分布式场景 & 层级设计复杂，需要记录阶段边界 \\
\bottomrule
\end{tabular}
\end{center}

\topic{工业参照：per-cpu counter 和 map-side combine}
Linux per-cpu counter 的核心思想，是高频写本 CPU 的局部状态，低频读取时再汇总。它牺牲了读取的即时一致性，换取写路径低争用。分布式计算中的 map-side combine 也是同一个思想：map worker 先本地聚合，把网络传输从“每条记录”降低到“每个本地 key 的 partial”。优秀设计不在名字，而在取舍：写快了，读或合并更复杂；状态分散了，观测和恢复要更清楚。

\topic{浮点归约的数值合同}
整数计数让 reduce 看起来简单，但很多工程计算是浮点。浮点归约不能只写“加法满足结合律”。有限精度下，合并顺序改变会改变舍入路径；FMA、fast-math、SIMD 水平合并和多线程树形合并都会影响结果。正确教材必须让读者在性能之前先选择数值合同。

\begin{center}
\begin{tabular}{p{0.22\linewidth}p{0.33\linewidth}p{0.33\linewidth}}
\toprule
合同 & 适用场景 & 代价 \\
\midrule
逐位一致 & 审计、回归测试、可复现构建 & 固定合并树，调度自由度下降 \\
误差范围 & 科学计算、统计、机器学习指标 & 需要误差测试和边界输入 \\
最快近似 & 临时监控、趋势估计 & 结果可能随线程数和机器变化 \\
补偿求和 & 高精度求和、抵消严重输入 & 计算更多，代码复杂 \\
\bottomrule
\end{tabular}
\end{center}

实验应构造四类输入：全正均匀数、极大极小混合、正负抵消、包含 NaN 或 Inf 的边界。比较串行左折叠、固定树形、动态树形、Kahan 或 pairwise sum。报告要同时给误差、耗时、重复运行一致性和编译选项。若只给“并行版快了”，没有说明误差，结论不完整。

\topic{Top-k 也是 reduce}
Top-k 合并可以看作 reduce：每个块产生局部候选集合，合并函数把两个候选集合合成新的前 k。它要求稳定 tie-break 和候选覆盖。若局部阶段提前丢掉全局仍需要的 key，后续 reduce 无法恢复。这个例子提醒读者：reduce 的合并函数不一定是加法，也可能是集合、堆、错误列表、统计摘要或 sketch。每一种都要写明 identity、合并顺序和信息丢弃边界。

\section{Scan：把未知输出数量变成确定位置}

Scan 也叫 prefix sum。算法题里它常被写成“计算前缀和”，但系统工程里更重要的作用是分配输出位置。稳定 filter、stream compaction、radix partition、稀疏矩阵构建、shuffle 文件 offset，都在回答同一个问题：每个并行任务要写到哪里。

\topic{从 \code{push_back} 推导 scan}
串行 filter 中，\code{push_back} 隐含维护了一个状态：前面已经输出了多少元素。并行后，这个状态不能由所有线程抢。稳定 filter 的三阶段推导如下：

\begin{enumerate}
  \item Count：每个 chunk 只统计自己会输出多少元素，不写全局输出。
  \item Exclusive scan：对 chunk counts 做前缀和，得到每个 chunk 的输出起点。
  \item Write：每个 chunk 重新扫描输入，把保留元素写入自己的连续区间。
\end{enumerate}

\begin{systemdiagram}{Stable filter：scan 把共享 push 改成独占区间}
\begin{pdfdiagram}
\node[book small node,text width=2.2cm] (c0) at (0,0) {chunk 0\\count = 2};
\node[book small node,text width=2.2cm] (c1) at (2.75,0) {chunk 1\\count = 3};
\node[book small node,text width=2.2cm] (c2) at (5.5,0) {chunk 2\\count = 1};

\node[book amber node,text width=2.2cm] (s0) at (0,-1.5) {start = 0};
\node[book amber node,text width=2.2cm] (s1) at (2.75,-1.5) {start = 2};
\node[book amber node,text width=2.2cm] (s2) at (5.5,-1.5) {start = 5};

\node[book teal node,text width=2.2cm] (w0) at (0,-3.0) {output[0..2)};
\node[book teal node,text width=2.2cm] (w1) at (2.75,-3.0) {output[2..5)};
\node[book teal node,text width=2.2cm] (w2) at (5.5,-3.0) {output[5..6)};

\node[book red node,text width=3.15cm] (claim) at (9.0,-1.5) {位置先算出来\\写出阶段不抢锁\\顺序由 chunk 顺序定义};

\draw[book strong arrow] (c0) -- (s0);
\draw[book strong arrow] (c1) -- (s1);
\draw[book strong arrow] (c2) -- (s2);
\draw[book strong arrow] (s0) -- (w0);
\draw[book strong arrow] (s1) -- (w1);
\draw[book strong arrow] (s2) -- (w2);
\draw[book dashed arrow] (s2) -- (claim);
\end{pdfdiagram}
\diagramnote{本图要证明的命题是：scan 的系统价值是把“不知道输出到哪里”改造成“每个 chunk 拥有一个确定输出区间”。}
\end{systemdiagram}

\topic{Inclusive 和 exclusive}
Inclusive scan 的第 \code{i} 个输出包含第 \code{i} 个输入；exclusive scan 的第 \code{i} 个输出只包含它之前的输入。分配输出起点通常使用 exclusive scan。API 名称必须写清楚，不要只叫 \code{scan}。

\begin{center}
\begin{tabular}{p{0.26\linewidth}p{0.24\linewidth}p{0.38\linewidth}}
\toprule
输入 & 输出 & 语义 \\
\midrule
\code{[3, 1, 4]} inclusive & \code{[3, 4, 8]} & 每个位置包含自身 \\
\code{[3, 1, 4]} exclusive & \code{[0, 3, 4]} & 每个位置是之前元素总量 \\
\bottomrule
\end{tabular}
\end{center}

\topic{Scan 的成本}
并行 scan 通常比串行 scan 做更多工作。局部扫描读输入写输出，扫描 block sums，回填 offset 还可能再读写输出。小输入上串行更快；大输入上并行才可能胜出。报告必须记录输入规模、block 数、临时数组大小和阶段耗时。Scan 是强工具，但不是免费工具。

\topic{错误边界}
Count 阶段和 write 阶段必须使用同一谓词版本。如果谓词依赖外部状态，count 得到的数量和 write 实际写出的数量可能不一致。工程实现应在 plan 或 manifest 中记录 predicate 版本、输入范围、expected count 和 checksum。数学上的 scan 正确，不代表系统里的 scan 一定正确。

\topic{块内位置和块间位置}
前面讲的是 chunk 级输出起点。若要让每个元素都知道自己的精确输出位置，还需要块内位置。常见做法是把谓词结果变成 0/1 mask，然后对块内 mask 做 exclusive scan。元素 \code{i} 的全局输出位置等于块起点加块内 rank。这个结构同时解释了 SIMD filter、GPU compaction 和数据库 selection vector 的位置生成。

\begin{lstlisting}[language=C++,caption={用 mask rank 表达元素输出位置}]
struct SelectedPosition {
  std::size_t input_index = 0;
  std::size_t output_index = 0;
};

SelectedPosition make_selected_position(std::size_t input_index,
                                        std::size_t block_output_begin,
                                        std::size_t rank_in_block) {
  return SelectedPosition{input_index, block_output_begin + rank_in_block};
}
\end{lstlisting}

这个示例故意只表达位置计算，不把谓词、scan 和写出混在一起。教学项目中可以把 mask、rank、output index 分别输出到报告，帮助读者定位错误：是谓词错、块内 scan 错、块间 offset 错，还是写出错。

\topic{Selection vector 和材料化输出}
Filter 后不一定立刻复制元素。若记录很大，可以输出 selection vector，只保存通过元素的下标；后续算子根据下标访问原始列。选择率低时，这能减少数据移动；选择率高时，下标本身也占带宽，后续间接访问可能破坏局部性。若后续要顺序扫描大量字段，材料化输出可能更好。工业执行器经常在 selection vector、bitmap 和 materialized batch 之间做取舍。

因此 scan 的结果可以服务多种输出形态：直接写紧凑数组、生成 selection vector、生成 bitmap rank、生成 block manifest。优秀实现不会把 scan 只绑定到一种容器，而是把“位置分配”作为可复用机制。

\section{Partition：多目的地的 scan}

Partition 把元素按 key、bucket 或目标 worker 分到多个输出区域。Filter 可以看成二路 partition：保留和不保留；多路 partition 则是每个 bucket 都要分配输出区间。它是单机并行算法和分布式 shuffle 之间最重要的桥。

\topic{二维计数表}
多路 partition 的第一步是局部直方图：每个 chunk 统计每个 bucket 有多少元素，形成 \code{counts[chunk][bucket]}。第二步对每个 bucket 沿 chunk 维度做 prefix sum，得到每个 chunk 在该 bucket 中的局部起点；同时计算每个 bucket 在全局输出中的基址。两者相加，才是最终写入位置。

\begin{systemdiagram}{多路 partition 的二维前缀}
\begin{pdfdiagram}
\node[book accent node,text width=2.5cm] (counts) at (0,0) {\code{counts}\\chunk x bucket};
\node[book teal node,text width=2.5cm] (bucket) at (3.35,0) {bucket totals\\全局基址};
\node[book amber node,text width=2.5cm] (prefix) at (6.7,0) {per-bucket scan\\chunk 偏移};
\node[book purple node,text width=2.5cm] (ranges) at (10.05,0) {output ranges\\manifest};

\draw[book strong arrow] (counts) -- (bucket);
\draw[book strong arrow] (bucket) -- (prefix);
\draw[book strong arrow] (prefix) -- (ranges);

\node[book small node,text width=2.6cm] at (0,-1.55) {局部形状};
\node[book small node,text width=2.6cm] at (3.35,-1.55) {每桶总大小};
\node[book small node,text width=2.6cm] at (6.7,-1.55) {每块写哪里};
\node[book small node,text width=2.6cm] at (10.05,-1.55) {可恢复身份};
\end{pdfdiagram}
\diagramnote{本图要证明的命题是：partition 是对每个 bucket 分别做 scan，manifest 是这些输出区间的工程化身份。}
\end{systemdiagram}

\topic{Manifest 是算法结果的一部分}
单机 partition 可以只返回一个重排后的数组，但工程系统还需要知道每个 bucket 的范围、记录数、字节数、checksum 和生成 attempt。这些信息就是 shuffle manifest 的雏形。

\begin{lstlisting}[language=C++,caption={PartitionManifestEntry 的教学形态}]
struct PartitionManifestEntry {
  std::int32_t partition_id = 0;
  std::uint64_t offset = 0;
  std::uint64_t records = 0;
  std::uint64_t bytes = 0;
  std::uint64_t checksum = 0;
};
\end{lstlisting}

Manifest 让失败恢复和审计成为可能。某个 chunk 写出失败时，系统知道哪些输出区间无效；重复执行同一个 chunk 时，系统知道它应该写回同一段；reduce 读取时，系统知道哪些 block 已提交。没有 manifest，partition 只是一次内存重排；有了 manifest，它就能成为分布式 shuffle 的可靠边界。

\topic{稳定和非稳定 partition}
稳定 partition 保持同一 bucket 内元素的原始相对顺序，适合审计、测试、可复现输出和需要稳定错误报告的场景。非稳定 partition 可以交换元素或按完成顺序写出，可能更快，但后续 stage 不能假设顺序。稳定性必须写进接口，不应由实现偶然决定。

\topic{热点 bucket}
Partition 解决写入位置，不自动解决负载均衡。若 70\% 记录落到同一个 bucket，下游 reduce 仍会长尾。工业系统会用局部 combine、热点 key 拆分、两级 hash、范围分区或采样来缓解。报告必须记录每个 bucket 大小和最大 bucket 占比，否则无法解释后续慢在哪里。

\topic{Radix partition：histogram、scan、scatter 的组合}
Radix partition 按 key 的若干 bit 分桶，常用于整数 key、哈希 join、radix sort 和数据库执行器。它的核心步骤正好是本章三件事的组合：先做局部 histogram，得到每个桶数量；再 scan 得到桶偏移；最后 scatter 到目标区间。每轮处理多少 bit 是重要参数：bit 多，桶多，counts 矩阵大，cache 压力高；bit 少，轮数多，读写次数增加。

这说明工业实现里的“参数调优”不是玄学。Radix bits 要根据 cache、TLB、输入规模、key 宽度和线程数测量。教材项目可以从 8 bit 一轮和 4 bit 一轮开始比较，记录 counts 矩阵大小、临时输出字节、总读写轮数和吞吐。读者会看到，算法复杂度相同，常数和内存布局仍能决定性能。

\topic{Partition 规则要版本化}
分区函数一旦改变，旧输出和新输出不能混用。分布式系统中，如果一部分 worker 使用旧 hash seed 或旧 bucket 数，另一部分使用新规则，同一个 key 会被送到不同 reducer。单机项目也应提前训练这个习惯：manifest 记录 partition version、bucket count、hash seed 或 range boundaries。Driver 发现版本不一致时拒绝合并。

\begin{lstlisting}[language=C++,caption={PartitionPlan 记录分区规则身份}]
struct PartitionPlan {
  std::uint32_t version = 0;
  std::uint32_t bucket_count = 0;
  std::uint64_t hash_seed = 0;
};
\end{lstlisting}

这个结构看起来简单，却能防止很多恢复事故。失败重试、增量执行、缓存复用和分布式 rolling upgrade 都需要知道输出是按哪套规则生成的。没有版本，manifest 只能说明“这里有一段数据”，不能说明“这段数据属于哪套分区语义”。

\topic{内存预算}
多路 partition 的额外内存来自 counts、offsets、临时输出和 manifest。若 \code{chunk_count} 很大、\code{bucket_count} 很大，二维 counts 可能很可观。报告应写出实际字节，而不是只写 \complexity{O(chunks * buckets)}。例如 counts 使用 \code{std::uint64_t}，空间就是 \code{chunk_count * bucket_count * 8} 字节；selection vector 若使用 \code{std::uint32_t}，最坏是 \code{input_count * 4} 字节。

内存预算会改变算法选择。桶数太多时，可以分批处理 bucket；热点明显时，可以先采样识别热点；输出太大时，可以按 batch 写临时 block；稳定性要求低时，可以使用局部 buffer 后合并。高质量工程不是只追求吞吐，而是在时间、空间、稳定性和恢复之间选择。

\section{工业界设计参照}

工业系统中的 reduce、scan 和 partition 很少以教科书名字出现，但它们的原理到处都是。优秀设计的共性，是把局部性、输出身份、失败恢复和确定性放在一起考虑。

\topic{Linux per-cpu counter}
per-cpu counter 把高频写放到每个 CPU 的局部槽位，读取全局值时再汇总。它的原理就是线程本地 partial。它牺牲了全局读的即时精确性，换取写路径低争用。迁移到本书项目，就是 worker 本地事件计数、任务完成数和 I/O 字节数；但若计数用于协议决策，例如“所有任务是否完成”，就不能使用近似汇总。

\topic{数据库 vectorized execution}
DuckDB、ClickHouse、Velox 等执行器常用 batch 和 selection vector。Filter 先生成 selection vector，而不是立刻复制整行；aggregate 先本地聚合，再合并；hash partition 先统计和分配区间，再把数据送到后续 stage。它们优秀的地方不是“用了向量化”四个字，而是让数据布局、输出身份和 pipeline breaker 都有明确边界。

\topic{Spark shuffle 和 MapReduce combine}
Map-side combine 是 reduce 的网络尺度版本：先本地聚合，再跨网络发送 partial。Shuffle manifest 是 partition 输出身份的工程版本：每个 block 有分区、大小、位置和校验。失败重试时，driver 通过 manifest 决定哪些 block 有效，哪些 attempt 应被丢弃。第 13 章的 partition manifest，正是这个工业设计的单机缩小版。

\topic{并行 STL 和 oneTBB}
并行算法库把范围切分、任务粒度、异常传播和合并策略封装起来。阅读它们时，不要只看是否开线程，而要看四件事：如何切 chunk，如何保存 partial，如何决定 cutoff，如何处理异常和取消。成熟库通常不会对所有输入强行并行，小输入会走串行或轻量路径。这是工业设计中的重要现实：并行也有固定成本。

\begin{sourcereadingbox}
工业案例阅读任务：选择 Linux per-cpu counter、DuckDB/ClickHouse aggregate、Spark shuffle、oneTBB parallel reduce 或 C++ 并行 STL 中一个实现。报告必须回答：它怎样切分输入；局部状态由谁拥有；合并顺序是否确定；输出身份如何记录；失败或异常如何处理；它和本章哪个实验可以对照。
\end{sourcereadingbox}

\section{确定性、审计和恢复}

并行算法常被默认允许“不确定顺序”，但工程项目不能这么粗糙。某些输出只关心集合内容，顺序无所谓；某些输出需要 bitwise reproducible，方便测试、缓存、审计和恢复；某些输出允许浮点误差，但要求误差范围稳定。确定性不是装饰，而是语义合同的一部分。

\topic{测试 oracle 由输出语义决定}
稳定 filter 可以逐项比较；非稳定输出要排序后比较或按多重集合比较；浮点 reduce 要按误差合同比较；partition 要比较每个 bucket 的内容和范围；manifest 要比较条目顺序、checksum 和版本。很多并行 bug 不是计算错，而是测试 oracle 和语义不匹配。

\topic{恢复边界}
Scan 产生的 offsets、partition 产生的 ranges、reduce 产生的 partial，都是恢复材料。若 write 阶段失败，可以根据 chunk 的输出区间重写；若 partition block 校验失败，可以只重做对应 chunk 或 bucket；若 dynamic reduce 结果不稳定，则恢复后可能无法 bitwise 对齐。算法结构直接影响恢复成本。

\topic{失败注入实验}
本章要求实现一个故障注入：count 和 scan 完成后，让某个 chunk 在 write 阶段失败。恢复程序应能知道该 chunk 的输出区间无效，其他 chunk 是否可保留由 manifest 决定。若输出由共享 \code{push_back} 产生，恢复程序很难定位边界。这个实验能让读者看到 scan 的系统价值，而不仅是性能价值。

\section{深度案例：把共享 push 改造成可恢复 stage}

现在把本章主线放进一个完整案例。输入是一批日志记录，目标是保留所有错误记录，并按错误码分区，供后续 reduce 统计错误类型。串行程序可以边扫边 \code{push_back} 到不同错误码 vector；并行后如果照搬这个结构，就会遇到锁、顺序、扩容和恢复问题。高质量设计要把它拆成四个 stage：count、scan、write、commit。

\topic{Stage 1：count 只收集形状}
每个 chunk 扫描输入，统计每个错误码会输出多少条记录。它不写最终输出，只产生 \code{counts[chunk][error_code]}。这个阶段容易重做，也容易校验。若 count 阶段已经看到非法格式，可以把错误写入 chunk-local error buffer，而不是立刻提交全局错误报告。

\topic{Stage 2：scan 生成地址合同}
对每个 error code 的 chunk counts 做 exclusive scan，得到每个 chunk 在该 error code 输出区域里的起点。这个结果就是地址合同：chunk 只能写自己的范围，不能越界，也不需要锁。若 write 阶段写出的数量和 count 不一致，说明谓词版本、输入或错误处理发生了变化，不能提交。

\topic{Stage 3：write 写入独占范围}
每个 chunk 重新扫描输入，把错误记录写到分配给自己的区间。为了保持稳定顺序，chunk 按输入范围顺序分配区间，chunk 内按原始顺序写。若业务不要求稳定顺序，可以选择更快的非稳定版本，但 manifest 必须标注输出顺序语义。

\topic{Stage 4：commit 提交 manifest}
write 完成后，系统为每个 error code 输出一个或多个 manifest 条目：错误码、offset、records、bytes、checksum、chunk id、attempt。只有 commit 成功的条目才被后续 stage 读取。若 chunk 失败，未提交条目被丢弃或标记失败；重试 chunk 时写回同一语义范围或新的 attempt 临时范围，再由 commit 选择有效版本。

\begin{systemdiagram}{共享 push 到可恢复 stage 的演化}
\begin{pdfdiagram}
\node[book red node,text width=2.25cm] (bad) at (0,0) {共享 vector\\push + lock};
\node[book accent node,text width=2.25cm] (count) at (2.75,0) {count\\局部形状};
\node[book amber node,text width=2.25cm] (scan) at (5.5,0) {scan\\地址合同};
\node[book teal node,text width=2.25cm] (write) at (8.25,0) {write\\独占区间};
\node[book purple node,text width=2.25cm] (commit) at (11.0,0) {commit\\manifest};

\draw[book strong arrow] (bad) -- (count);
\draw[book strong arrow] (count) -- (scan);
\draw[book strong arrow] (scan) -- (write);
\draw[book strong arrow] (write) -- (commit);

\node[book label,text width=10.8cm] at (5.5,-1.35) {重构后的 stage 不只是更快，还让输出位置、顺序、校验和失败重做都有明确边界。};
\end{pdfdiagram}
\diagramnote{本图要证明的命题是：scan 把并发写输出的性能问题，同时转化成可恢复的提交协议问题。}
\end{systemdiagram}

\topic{为什么这是工业级思路}
这个案例的四阶段结构和数据库执行器、shuffle 写出、日志处理管线非常接近。工业系统常把正常路径和恢复路径一起设计：正常路径要少锁、顺序写、批量处理；恢复路径要知道哪个输入产生哪个输出，重复执行是否幂等，哪些输出已经提交。共享 \code{push_back} 的最大问题不是慢，而是它没有留下这些边界。

\section{项目落地}

Compute Systems Engine 在本章应新增一个小型并行算法套件。它不需要一次写成通用框架，但必须覆盖 reduce、scan 和 partition 的最小可靠切片。

\begin{center}
\begin{tabular}{p{0.2\linewidth}p{0.35\linewidth}p{0.33\linewidth}}
\toprule
模块 & 最小能力 & 必交证据 \\
\midrule
parallel count-if & 串行 reference、global atomic 反例、partial 版本 & 正确性、线程扩展、共享写解释 \\
stable filter & count、exclusive scan、write 三阶段 & 输出顺序、chunk counts、offsets、故障注入 \\
multiway partition & 二维 counts、per-bucket scan、manifest & 每桶大小、最大倾斜、checksum、稳定性 \\
deterministic reduce & 固定合并树、误差合同 & 浮点误差、重复运行一致性、性能代价 \\
\bottomrule
\end{tabular}
\end{center}

\topic{输入族}
输入不能只有随机均匀数组。必须覆盖空输入、一个元素、刚好一个 chunk、最后 chunk 不满、全通过、全不通过、低选择率、高选择率、单热点 bucket、Zipf 倾斜、浮点极大极小混合和失败 chunk。压力输入要攻击合同边界，而不是只扩大规模。

\topic{Cutoff}
并行算法不是输入越小越应该并行。任务提交、队列调度、scan、partial、合并都有固定成本。项目应测量串行和并行交叉点，形成 cutoff。小输入走串行，大输入走并行；中等输入可能单线程 SIMD 更合适。Cutoff 不是永恒常数，而是随机器、编译器和输入分布变化的配置。

\topic{报告模板}
本章实验报告固定写八段：串行语义；并行合同；切分和合并策略；输入族；阶段账本；正确性 oracle；性能结果；未验证边界。阶段账本至少包含 input records、output records、chunk count、temporary bytes、partial bytes、max bucket、checksum 和 elapsed time。

\begin{center}
\begin{tabular}{p{0.19\linewidth}p{0.34\linewidth}p{0.35\linewidth}}
\toprule
实验 & 必测输入 & 必交结论 \\
\midrule
parallel count-if & 空输入、全命中、全不命中、随机命中、小输入、大输入 & 全局 atomic 为什么慢，partial 何时值得，cutoff 在哪里 \\
浮点 reduce & 极大极小混合、正负抵消、NaN/Inf、重复运行 & 误差合同、固定树形代价、fast 模式边界 \\
stable filter & 0\%、1\%、50\%、100\% 选择率，最后 chunk 不满 & scan 如何分配位置，稳定顺序成本，故障注入结果 \\
multiway partition & 均匀 key、单热点、Zipf 倾斜、bucket 数变化 & counts 矩阵大小、最大 bucket、manifest 是否足够恢复 \\
radix partition & 4 bit、8 bit、不同 key 宽度和输入规模 & 桶数、轮数、cache 压力和吞吐之间的取舍 \\
\bottomrule
\end{tabular}
\end{center}

\topic{阶段账本}
每个实验都要输出阶段账本，而不是只输出总耗时。建议字段如下：

\begin{lstlisting}[caption={并行算法阶段账本字段}]
stage,input_records,output_records,chunks,buckets,temporary_bytes,
partial_bytes,max_bucket,checksum,elapsed_ns,mode
\end{lstlisting}

这些字段帮助读者解释现象。若并行 filter 在 1\% 选择率下不快，可能是 count 和 scan 固定成本主导；若 partition 在单热点输入下慢，\code{max_bucket} 会暴露下游长尾；若 radix partition 的 8 bit 版本变慢，\code{buckets} 和 \code{temporary_bytes} 能说明 cache 压力。没有阶段账本，性能解释很容易变成猜测。

\begin{labbox}
实验：把共享 \code{push_back} filter 改造成 stable filter。实现串行 reference、加锁 push 反例、每线程本地 buffer 版本、count + scan + write 版本。比较输出顺序、锁等待或替代指标、临时内存、不同选择率下的耗时，并加入一个 write 阶段故障注入。
\end{labbox}

\begin{rubricbox}
本章大作业按 100 分评分。正确性 30 分：reference、oracle、边界输入和重复运行一致。原理推导 20 分：能从串行语义推导 partial、scan offsets 和 partition ranges。实验质量 20 分：覆盖选择率、热点、浮点误差、cutoff 和故障注入。工业迁移 15 分：能对照 per-cpu、vectorized execution、shuffle 或并行库设计讲清取舍。工程表达 15 分：代码接口位宽明确，manifest 和阶段报告可复查。
\end{rubricbox}

\section{复查清单}

\topic{一、Reduce 是否写清合并语义}
检查 identity、结合性、交换性、确定性、浮点误差和 partial 所有权。若这些没有写清，就不能称为可靠并行 reduce。

\topic{二、Scan 是否解释输出位置}
检查读者是否能从共享 \code{push_back} 推导出 count、exclusive scan、write 三阶段。若只会背 prefix sum 公式，说明还没学到系统意义。

\topic{三、Partition 是否留下 manifest}
检查每个 bucket 的范围、记录数、字节数和 checksum 是否可见。没有 manifest，后续 shuffle 和恢复会断链。

\topic{四、是否区分稳定和非稳定输出}
稳定性影响测试、审计、缓存和幂等。接口和报告必须说明是否保持顺序，以及为此付出了什么成本。

\topic{五、是否记录资源预算}
Partial、counts、offsets、本地 buffer、manifest 都占内存。报告要写实际字节，不只写复杂度。

\topic{六、是否连接下一章运行时}
Reduce、scan 和 partition 都形成任务依赖。局部任务完成后才能 scan，scan 完成后才能 write，所有 partition block 提交后才能 reduce。下一章任务运行时要把这些依赖显式调度，而不是让 worker 阻塞等待。

\section{本章和任务运行时的连接}

Reduce、scan 和 partition 一旦写成 stage，就自然产生任务图。Count tasks 可以并行；scan offsets 必须等所有 count 完成；write tasks 必须等 offsets 完成；reduce tasks 必须等对应 partition block 提交。若用普通线程池，最容易写成 worker 内部阻塞等待下一阶段，造成线程池耗尽或低效 barrier。下一章的任务运行时要解决的，正是如何把这些依赖显式化，让 worker 只执行 ready task，而不是占着线程等待。

\topic{从算法阶段到任务图节点}
一个 stable filter 可以拆成三类节点：count node、scan node、write node。Count nodes 只依赖输入 chunk；scan node 依赖所有 count nodes；write nodes 依赖 scan node 和对应 input chunk。Partition 多一个 bucket 维度；reduce 多一个合并树维度。这样算法结构直接变成运行时调度结构。

\begin{systemdiagram}{Stable filter 的任务依赖}
\begin{pdfdiagram}
\node[book small node,text width=2.1cm] (c0) at (0,0) {count 0};
\node[book small node,text width=2.1cm] (c1) at (2.6,0) {count 1};
\node[book small node,text width=2.1cm] (c2) at (5.2,0) {count 2};
\node[book amber node,text width=2.4cm] (scan) at (2.6,-1.55) {scan offsets};
\node[book teal node,text width=2.1cm] (w0) at (0,-3.1) {write 0};
\node[book teal node,text width=2.1cm] (w1) at (2.6,-3.1) {write 1};
\node[book teal node,text width=2.1cm] (w2) at (5.2,-3.1) {write 2};
\node[book purple node,text width=2.5cm] (commit) at (8.4,-3.1) {commit manifest};

\draw[book strong arrow] (c0) -- (scan);
\draw[book strong arrow] (c1) -- (scan);
\draw[book strong arrow] (c2) -- (scan);
\draw[book strong arrow] (scan) -- (w0);
\draw[book strong arrow] (scan) -- (w1);
\draw[book strong arrow] (scan) -- (w2);
\draw[book strong arrow] (w0) -- (commit);
\draw[book strong arrow] (w1) -- (commit);
\draw[book strong arrow] (w2) -- (commit);
\end{pdfdiagram}
\diagramnote{本图要证明的命题是：并行算法的阶段依赖应该交给任务运行时调度，而不是让 worker 线程阻塞等待。}
\end{systemdiagram}

\topic{运行时要保存哪些指标}
为了让下一章能分析任务运行时，本章的实验应该输出每个 task 的 submit、ready、start、finish 时间，以及输入记录数、输出记录数和错误数。这样第 14 章就能解释 worker 空闲、队列长尾、barrier 等待和任务粒度。若第 13 章只输出总耗时，第 14 章就没有诊断材料。

\topic{从 barrier 到依赖调度}
最简单实现是在每个阶段之间放 barrier：所有 count 完成后主线程做 scan，再启动所有 write。这个版本容易理解，适合教学起点。更进一步，运行时可以把依赖显式化：当最后一个 count 完成时，scan node 变 ready；scan 完成后，所有 write nodes 变 ready；write 完成后，commit node 变 ready。这种结构减少不必要等待，也为 work stealing、故障重试和背压打基础。

\section{本章习题}

\begin{exercisebox}
习题一：为 parallel count-if 写完整正确性证明。证明每个输入元素被统计一次且只统计一次，partial 合并结果等于串行 reference。

习题二：设计 inclusive scan 和 exclusive scan 的 API。给出输入 \code{[3, 1, 4]} 的两个输出，并说明哪个更适合分配输出偏移。

习题三：为 stable filter 写三阶段伪代码。要求说明每阶段读写哪些数组、是否并行、是否有共享写、需要哪些临时空间。

习题四：设计一个倾斜 key 的 partition 测试。要求解释为什么写入区间没有冲突，但下游 reduce 仍可能负载不均。

习题五：比较线性合并和树形合并。给出 partial 数为 8 时的合并步骤，并说明什么时候树形合并不值得。

习题六：阅读一个工业案例，例如 per-cpu counter、vectorized aggregate 或 Spark shuffle manifest。写出它的真实约束、核心原理、取舍、和本章实验的对应关系。
\end{exercisebox}
